\chapter{Transient Ischaemic Attack (TIA)}
\index{Transient Ischaemic Attack (TIA)}
\medskip
\noindent\textit{Educational reference; always seek professional advice for individual care.}

\section*{Definition and Overview}
A transient ischaemic attack (TIA) is a brief episode of neurological dysfunction caused by temporary reduction of blood flow to the brain, spinal cord, or retina, without lasting tissue infarction on appropriate imaging. Symptoms resemble stroke and may include sudden weakness, facial drooping, speech difficulty, loss of vision in one eye, or loss of balance. Symptoms often resolve within minutes, but a TIA is still a medical emergency because it can precede a completed stroke.

\section*{Recognising a TIA}
Symptoms usually start suddenly and are maximal at onset. Possible features include:
\begin{itemize}
    \item weakness or numbness, especially on one side of the face or body
    \item difficulty speaking, understanding speech, or finding words
    \item sudden loss or blurring of vision in one or both eyes
    \item sudden dizziness, loss of coordination, or difficulty walking
    \item rarely, brief double vision or loss of part of the visual field
\end{itemize}

Use the FAST warning signs: Face drooping, Arm weakness, Speech difficulty, and Time to seek emergency help. A symptom that has completely resolved still needs urgent assessment. Do not drive yourself if stroke-like symptoms are occurring.

\section*{Causes and Risk Factors}
A TIA may result from a small clot travelling from the heart or a large artery, narrowing of a carotid artery, or disease of small brain arteries. Important risk factors include high blood pressure, atrial fibrillation, diabetes, smoking, high cholesterol, sleep apnoea, obesity, older age, and previous TIA or stroke. Some people have no obvious risk factor.

\section*{When to Seek Medical Care}
\textbf{Contact local emergency services immediately} for any sudden facial droop, weakness, numbness, speech difficulty, vision loss, severe imbalance, or other new stroke-like symptom, even if it improves. Note the time the person was last known to be well. Keep the person safe, do not give food, drink, or medicines unless advised, and follow the emergency operator's instructions.

\section*{Diagnosis and Investigations}
Urgent assessment may include a neurological examination, blood-pressure measurement, blood tests, glucose testing, ECG and prolonged rhythm monitoring, brain imaging (usually CT or MRI), and imaging of the carotid and other relevant arteries. The clinician will look for conditions that mimic TIA, such as migraine aura, seizure, low blood glucose, fainting, vestibular disorders, or functional neurological symptoms. A normal early scan does not make a new sudden neurological symptom safe to ignore.

\section*{Management}
Treatment depends on the cause and bleeding risk. It may include antiplatelet therapy, anticoagulation for atrial fibrillation, a statin, blood-pressure and diabetes management, smoking cessation, and treatment of sleep apnoea. Severe carotid narrowing may require a vascular procedure. Medicines should be started, stopped, or combined only under clinical guidance because inappropriate antithrombotic treatment can cause serious bleeding.

\section*{Prevention and Follow-up}
Attend the recommended rapid-access or specialist review, take medicines as prescribed, and keep follow-up appointments. A heart-healthy eating pattern, regular activity appropriate to ability, tobacco cessation, moderation of alcohol, and management of blood pressure, cholesterol, diabetes, and weight reduce future risk. Family members should learn the local emergency warning signs and response number.

\section*{Outlook}
The risk of stroke is highest soon after a TIA, particularly when symptoms are recent, recurrent, or caused by atrial fibrillation or carotid disease. Prompt assessment and targeted prevention substantially reduce that risk. A TIA does not guarantee that a stroke will occur, but it is a warning that needs action.

\section*{Sources and Further Reading}
\begin{itemize}
    \item World Health Organization. \textit{Stroke.} https://www.who.int/health-topics/stroke
    \item American Stroke Association. \textit{Transient ischemic attack (TIA).} https://www.stroke.org/
    \item NHS. \textit{Transient ischaemic attack (TIA).} https://www.nhs.uk/conditions/transient-ischaemic-attack-tia/
    \item National Institute of Neurological Disorders and Stroke. \textit{Stroke information.} https://www.ninds.nih.gov/health-information/disorders/stroke
\end{itemize}
